% RESULTADO_RECUPERADO: funtor dimensional integral, eq:cZ-internos,
% eq:longitud-ruta; composición reunida en Artículo V REV02, sección 69.
% Incorporación autónoma a III: mismo transporte, mismas hojas y misma carta.
\subsection{Vitesse constitutive et réalisation cinématique de l’horloge}
\label{iii:sec:enlace-velocidad-constitutiva}

La réponse constitutive et le lecteur cinématique réalisent une même
section de vitesse. Leur identification conserve l’ordre de construction :
les feuillets d’APP, leur orientation TRIT et le transport angulaire exprimé
par le TPK précèdent les lectures de vitesse et le choix des unités.
Dans la représentation à deux feuillets de la
section~\ref{iii:sec:hojas}, avec $x>|y|$ et les projecteurs de rang un
$P_\pm$, le transport est
\[
 T=e^{-xI-y\mathcal R}=q_+P_++q_-P_-,
 \qquad q_\pm=e^{-x\mp y},\qquad 0<q_\pm<1.
\]
Ses blocs temporels de longueurs $90$ et $120$ produisent la réponse
$\mathcal V$ de \eqref{iii:eq:vacuum-op}. Les coordonnées $x,y$ et les
canaux sont des sorties du développement angulaire préalable ; cette composition
applique ses lecteurs à ce transport déjà construit.

Le lecteur logarithmique symétrique du transport temporel est
\begin{equation}
 \mathcal E=
 \log\!\bigl[(1-q_+^{90})(1-q_-^{90})\bigr]
 -\log\!\bigl[(1-q_+^{120})(1-q_-^{120})\bigr].
 \label{iii:eq:velocidad-lector-simetrico}
\end{equation}
Sur la droite dimensionnelle de vitesse, avec une section positive $u_C$, le
lecteur cinématique exprime $c_{\mathrm{int}}=\exp(-\mathcal E)u_C$.
La réponse constitutive exprime
$c_{\mathrm{vac}}=\widehat c u_C$ au moyen de \eqref{iii:eq:four}.

\begin{proposition}[Identité constitutive et cinématique]
\label{iii:prop:identidad-velocidad}
Dans la même coordinatisation dimensionnelle, on vérifie
\begin{equation}
 c_{\mathrm{int}}=c_{\mathrm{vac}}=\widehat c\,u_C,
 \qquad \widehat c=(r_+r_-)^{-1}=\exp(-\mathcal E).
 \label{iii:eq:identidad-velocidad}
\end{equation}
L'identification conserve l'échange des feuillets.
\end{proposition}
\begin{proof}
Tous les facteurs de \eqref{iii:eq:velocidad-lector-simetrico} sont
positifs. La loi du logarithme et \eqref{iii:eq:rchannels} donnent
\[
 \mathcal E=\log(r_+r_-)=\log\det\mathcal V,
 \qquad e^{-\mathcal E}=(\det\mathcal V)^{-1}.
\]
Le déterminant correspond à la réalisation à deux feuillets de rang un.
Les deux définitions expriment donc la même section. L’échange
de feuillets permute $r_+$ et $r_-$ et conserve leur produit. Le déterminant
$\det T=e^{-2x}$ correspond au transport élémentaire, tandis que
$\det\mathcal V=r_+r_-$ correspond à sa réponse temporelle ;
\eqref{iii:eq:identidad-velocidad} utilise le second.
\end{proof}

\subsubsection{Longueur de trajectoire et durée élémentaire}

Le repère dimensionnel positif $(u_S,u_C,u_T,u_Q)$ induit la base de longueur $u_L=u_Cu_T$. Pour une trajectoire enrichie $\gamma$ de $n=|\gamma|$ transitions et à canal constant, les lecteurs additifs sont
\begin{equation}
 T_{\mathrm{int}}(\gamma)=nu_T,
 \qquad L_{\mathrm{int}}(\gamma)=n\widehat c\,u_L.
 \label{iii:eq:velocidad-longitud-reloj}
\end{equation}
La concaténation additionne le nombre de transitions et leurs réalisations.
Les routes enrichies conservent en outre leur orientation et leur mémoire.

\begin{corollary}[Normalisation de la longueur élémentaire]
Pour $n>0$, $L_{\mathrm{int}}(\gamma)/T_{\mathrm{int}}(\gamma)
=c_{\mathrm{vac}}$. Dans la normalisation $t_0=u_T$,
\begin{equation}
 \ell_0=c_{\mathrm{int}}t_0=\widehat c\,u_L.
 \label{iii:eq:longitud-elemental-constitutiva}
\end{equation}
Dans une coordinatisation temporelle $t_0=\tau_0u_T$, avec $\tau_0>0$, on obtient
$\ell_0=\tau_0\widehat c\,u_L$.
\end{corollary}
\begin{proof}
Diviser les deux expressions de \eqref{iii:eq:velocidad-longitud-reloj} simplifie $n$ et utilise $u_L/u_T=u_C$. Multiplier la section de vitesse par $t_0$ prouve les deux formules élémentaires.
\end{proof}

Lorsque le canal dépend de l’arête, le lecteur conserve la forme additive
$L_{\mathrm{int}}(\gamma)=\sum_{e\in\gamma}\widehat c(e)u_L$.
Son quotient par $nu_T$ est la moyenne des vitesses de ces arêtes.
L’expression à canal constant est la spécialisation homogène de
cette lecture.

\subsubsection{Changement de coordinatisation et normalisation adaptée}

\begin{proposition}[Covariance de l’identification]
Soient $u_C'=d u_C$ et $u_T'=\eta u_T$, avec $d,\eta>0$. Les coordonnées des mêmes sections satisfont
\begin{equation}
 \widehat c'=d^{-1}\widehat c,
 \qquad \tau_0'=\eta^{-1}\tau_0,
 \qquad u_L'=d\eta u_L.
 \label{iii:eq:velocidad-cambio-carta}
\end{equation}
Si une autre représentation écrit $c'=\widehat c' u_C'$ avec $u_C'=d u_C$,
l’égalité des sections équivaut exactement à $d\widehat c'=\widehat c$.
\end{proposition}
\begin{proof}
L'égalité de chaque section dans les deux bases donne $\widehat c'u_C'=\widehat c u_C$ et $\tau_0'u_T'=\tau_0u_T$. En particulier,
\[
 \tau_0'\widehat c'u_L'
 =\frac{\tau_0}{\eta}\frac{\widehat c}{d}\,d\eta u_L
 =\tau_0\widehat c u_L=\ell_0.
\]
La dernière équivalence provient de la même substitution pour la vitesse.
\end{proof}

Le choix adapté $d=\widehat c$ donne $u_C'=c_{\mathrm{vac}}$ et
$\widehat c'=1$ : il exprime la section construite dans une base adaptée.
Le coefficient de la coordinatisation interne reste $(r_+r_-)^{-1}$, et c’est
la normalisation utilisée par l’inversion cubique. En maintenant les
bases d’action et de charge, \eqref{iii:eq:dimensions} produit
\[
 \widehat\mu'=d\widehat\mu,\qquad
 \widehat\varepsilon'=d\widehat\varepsilon,\qquad
 \widehat\mu'\widehat\varepsilon'=(\widehat c')^{-2}.
\]
Dans la coordinatisation adaptée, on a $\widehat\mu'=r_-/r_+$ et
$\widehat\varepsilon'=r_+/r_-$, par substitution de
$\widehat\mu=r_-^2$ et $\widehat\varepsilon=r_+^2$.

\subsubsection{Conservation du lecteur sous amplification}

Pour le raffinement $\mathcal V_m=\mathcal V\otimes I_{9^m}$, la trace normalisée $\tau_m=\operatorname{Tr}/(2\cdot9^m)$ satisfait
\begin{equation}
 \exp[-2\tau_m(\log\mathcal V_m)]
 =\exp[-\log r_+-\log r_-]=\widehat c.
 \label{iii:eq:velocidad-traza-normalizada}
\end{equation}
En effet, chaque valeur propre apparaît $9^m$ fois ; cette multiplicité
s’annule avec le dénominateur de la trace normalisée. Le déterminant
ordinaire amplifié est $(r_+r_-)^{9^m}$. La normalisation à deux feuillets,
exprimée de manière équivalente par \eqref{iii:eq:velocidad-traza-normalizada},
conserve le lecteur constitutif pendant l’amplification. Cette
réalisation conserve le transport enrichi dont elle est issue.
